Conventions: the chart's complex unit is $I_{6}$, the complex unit of the conformal spinor space $\CS$, acting on conformal spinors by right multiplication; numbers formed from conformal spinors (values of the Hermitian form, components, traces) are read in $\mathrm{span}\{1, I_{6}\}$, and the bar is the conjugation that reverses $I_{6}$; $\Som$ and $\Spi$ are the halves of $\CS$ spanned by $e_{1}, e_{2}$ and by $e_{3}, e_{4}$, the $E = -1$ and $E = +1$ eigenspaces.

(1) Expand $Z^{1}\wedge Z^{2} = (\omega^{1} + \pi^{1})\wedge(\omega^{2} + \pi^{2})$ into four terms. Pairing against $n_{\infty} = e_{1}\wedge e_{2}$ annihilates every term with an $\omega$ factor (the wedge repeats a vector of $\Som$), leaving $\varepsilon(e_{1}\wedge e_{2}\wedge\pi^{1}\wedge\pi^{2}) = \varepsilon_{a'b'}\pi_{1}^{a'}\pi_{2}^{b'}\,\varepsilon(e_{1}\wedge e_{2}\wedge e_{3}\wedge e_{4}) = \br{\pi^{1}\pi^{2}}$; pairing against $n_{0} = e_{3}\wedge e_{4}$ likewise leaves $\varepsilon(e_{3}\wedge e_{4}\wedge\omega^{1}\wedge\omega^{2}) = [\omega^{1}\omega^{2}]\,\varepsilon(e_{3}\wedge e_{4}\wedge e_{1}\wedge e_{2}) = [\omega^{1}\omega^{2}]$, the permutation being even. Linearity in $I_{\Lambda}$ gives \eqref{eq:blambda}; at $\Lambda = 0$ the frozen point is $n_{\infty}$ (the infinity twistor) (Lemma \claimref{mass.de-sitter-scale-frozen-point}(3)).

(2) \emph{Slots of the grip.} $M_{U\wedge I}$ is linear in $U \wedge I$, so with $I_{\Lambda} = n_{\infty} + \lambda n_{0}$ the grip is $P = P_{\infty} + \lambda P_{0}$, where $P_{\infty} := M(n_{\infty})$ and $P_{0} := M(n_{0})$. Put $d := M_{n_{\infty}\wedge n_{0}}$, and for $x$ in the $x$-slot $p(x) := M_{x\wedge n_{\infty}}$, $k(x) := M_{x\wedge n_{0}}$, three linear forms in $x$. Then $\ip{P}{n_{\infty}} = M_{n_{\infty}\wedge I_{\Lambda}} = \lambda M_{n_{\infty}\wedge n_{0}} = \lambda d$, $\ip{P}{n_{0}} = M_{n_{0}\wedge n_{\infty}} = -d$, and $\ip{P}{x} = p(x) + \lambda k(x)$. Writing $P = \alpha n_{\infty} + \beta n_{0} + v$ with $v$ in the $x$-slot, $\alpha = \ip{P}{n_{0}} = -d$, $\beta = \ip{P}{n_{\infty}} = \lambda d$, and $\ip{v}{x} = p(x) + \lambda k(x)$ for all $x$. Hence
\[
Q(P) = 2\alpha\beta + Q(v) = -2\lambda d^{2} + Q(p) + 2\lambda\, p\cdot k + \lambda^{2}\, Q(k),
\]
where $f \cdot g$ denotes the dual metric on linear forms on the $x$-slot. By Remark \ref{rem:nullity} the metric matrix on the $x$-slot, in the components $x^{aa'}$, is $G = -\varepsilon\otimes\varepsilon$; since $\varepsilon^{2} = -1$, $G^{2} = (\varepsilon\otimes\varepsilon)^{2} = \varepsilon^{2}\otimes\varepsilon^{2} = 1$, so $G^{-1} = G$ and $f\cdot g = -\varepsilon_{ab}\varepsilon_{a'b'}f_{aa'}g_{bb'}$ with the same numerical symbols. Below, $u\cdot v$ for two $2$-component objects is the plain sum $\sum_{a}u_{a}v_{a}$ of products of components.

\emph{The three forms in components.} By Lemma \claimref{mass.chart-form-spin-isomorphism}(2), $A(x\wedge n_{\infty})Z = b_{x}\pi \in \Som$, so by \eqref{eq:hermitian} and \eqref{eq:moment}, $p(x) = I_{6}\sum_{i}\hp{\pi^{i}}{b_{x}\pi^{i}} = I_{6}\sum_{i} x^{aa'}\,\overline{\pi^{i}_{a}}\,(\varepsilon\pi^{i})_{a'}$, where $(\varepsilon\pi)_{a'} := \varepsilon_{a'b'}\pi^{b'}$ and $\overline{\pi^{i}_{a}}$ is the conjugate of the component of $\pi^{i}$ paired with $e_{a}$ by $\Sigma$. Thus $p = \sum_{i}p^{i}$ with $p^{i}_{aa'} = I_{6}\,\overline{\pi^{i}_{a}}\,(\varepsilon\pi^{i})_{a'}$. By Lemma \claimref{mass.chart-form-spin-isomorphism}(3), $A(x\wedge n_{0})Z = c_{x}\omega \in \Spi$, so $k(x) = I_{6}\sum_{i}\overline{\hp{c_{x}\omega^{i}}{\omega^{i}}} = I_{6}\sum_{i}\sum_{a'}(c_{x}\omega^{i})^{a'}\,\overline{\omega^{i}_{a'}} = -I_{6}\sum_{i}x^{ba'}\varepsilon_{ba}\omega^{i\,a}\,\overline{\omega^{i}_{a'}}$; thus $k = \sum_{i}k^{i}$ with $k^{i}_{aa'} = -I_{6}\,(\varepsilon\omega^{i})_{a}\,\overline{\omega^{i}_{a'}}$, where $(\varepsilon\omega)_{a} := \varepsilon_{ab}\omega^{b}$. By Lemma \claimref{mass.chart-form-spin-isomorphism}(1), $A(n_{\infty}\wedge n_{0})Z = (-\tfrac12\omega, \tfrac12\pi)$, so $\bar{Z}\cdot AZ = -\tfrac12\hp{\pi}{\omega} + \tfrac12\overline{\hp{\pi}{\omega}} = -I_{6}\,\mathrm{Im}\hp{\pi}{\omega}$ and $d = \sum_{i}\mathrm{Im}\,t_{ii} = \mathrm{Im}\tr t$.

\emph{The three pairings.} Using $\varepsilon^{\mathsf T}\varepsilon = 1$, $\varepsilon^{2} = -1$ and $\det\varepsilon = 1$ for the $2\times 2$ symbol:
\begin{align*}
p^{i}\cdot p^{l} &= -\varepsilon_{ab}\varepsilon_{a'b'}\,(I_{6})^{2}\,\overline{\pi^{i}_{a}}\,\overline{\pi^{l}_{b}}\,(\varepsilon\pi^{i})_{a'}(\varepsilon\pi^{l})_{b'}
= \overline{\br{\pi^{i}\pi^{l}}}\;\br{\varepsilon\pi^{i},\varepsilon\pi^{l}} = |\br{\pi^{i}\pi^{l}}|^{2},\\
k^{i}\cdot k^{l} &= -\varepsilon_{ab}\varepsilon_{a'b'}\,(-I_{6})^{2}\,(\varepsilon\omega^{i})_{a}(\varepsilon\omega^{l})_{b}\,\overline{\omega^{i}_{a'}}\,\overline{\omega^{l}_{b'}}
= [\varepsilon\omega^{i},\varepsilon\omega^{l}]\;\overline{[\omega^{i}\omega^{l}]} = |[\omega^{i}\omega^{l}]|^{2},\\
p^{i}\cdot k^{l} &= -\varepsilon_{ab}\varepsilon_{a'b'}\,(I_{6})(-I_{6})\,\overline{\pi^{i}_{a}}\,(\varepsilon\omega^{l})_{b}\,(\varepsilon\pi^{i})_{a'}\,\overline{\omega^{l}_{b'}}
= -\bigl(-\overline{\pi^{i}}\cdot\omega^{l}\bigr)\bigl(\pi^{i}\cdot\overline{\omega^{l}}\bigr) = t_{il}\,\overline{t_{il}} = |t_{il}|^{2},
\end{align*}
where in the last line $\varepsilon_{ab}u_{a}(\varepsilon v)_{b} = u^{\mathsf T}\varepsilon\varepsilon v = -u\cdot v$ and $\varepsilon_{a'b'}(\varepsilon u)_{a'}v_{b'} = u^{\mathsf T}\varepsilon^{\mathsf T}\varepsilon v = u\cdot v$. Summing over $i, l$: $Q(p) = 2|\br{\pi^{1}\pi^{2}}|^{2}$ (the diagonal terms vanish, and this is Proposition \claimref{mass.mass-symplectic-area-two-twistor-formula} again), $Q(k) = 2|[\omega^{1}\omega^{2}]|^{2}$, and $p\cdot k = \sum_{i,l}|t_{il}|^{2}$. Therefore
\[
Q(P) = 2|\br{\pi^{1}\pi^{2}}|^{2} + 2\lambda^{2}|[\omega^{1}\omega^{2}]|^{2} + 2\lambda\Bigl(\sum_{i,l}|t_{il}|^{2} - (\mathrm{Im}\tr t)^{2}\Bigr).
\]

\emph{The comparison.} $2|B_{\Lambda}|^{2} = 2|\br{\pi^{1}\pi^{2}}|^{2} + 2\lambda^{2}|[\omega^{1}\omega^{2}]|^{2} + 4\lambda\,\mathrm{Re}\bigl(\br{\pi^{1}\pi^{2}}\,\overline{[\omega^{1}\omega^{2}]}\bigr)$. The matrix $t$ is the product of the matrix with rows $\overline{\pi^{1}}, \overline{\pi^{2}}$ and the matrix with columns $\omega^{1}, \omega^{2}$, so $\det t = \overline{\br{\pi^{1}\pi^{2}}}\,[\omega^{1}\omega^{2}]$ and $\mathrm{Re}\det t = \mathrm{Re}\bigl(\br{\pi^{1}\pi^{2}}\overline{[\omega^{1}\omega^{2}]}\bigr)$. It remains to show
\[
\sum_{i,l}|t_{il}|^{2} - (\mathrm{Im}\tr t)^{2} - 2\,\mathrm{Re}\det t \;=\; C .
\]
With $h_{11} = 2\,\mathrm{Re}\,t_{11}$, $h_{22} = 2\,\mathrm{Re}\,t_{22}$, $h_{12} = t_{12} + \overline{t_{21}}$, the right side is $(\mathrm{Re}\,t_{11} - \mathrm{Re}\,t_{22})^{2} + |t_{12}|^{2} + |t_{21}|^{2} + 2\,\mathrm{Re}(t_{12}t_{21})$, and the left side is $|t_{11}|^{2} + |t_{22}|^{2} + |t_{12}|^{2} + |t_{21}|^{2} - (\mathrm{Im}\,t_{11} + \mathrm{Im}\,t_{22})^{2} - 2\,\mathrm{Re}(t_{11}t_{22}) + 2\,\mathrm{Re}(t_{12}t_{21})$. Their difference is
\[
|t_{11}|^{2} + |t_{22}|^{2} - (\mathrm{Im}\,t_{11} + \mathrm{Im}\,t_{22})^{2} - 2\,\mathrm{Re}(t_{11}t_{22}) - (\mathrm{Re}\,t_{11} - \mathrm{Re}\,t_{22})^{2},
\]
and writing $t_{jj} = r_{j} + I_{6} s_{j}$ this is $(r_{1}^{2} + s_{1}^{2} + r_{2}^{2} + s_{2}^{2}) - (s_{1} + s_{2})^{2} - 2(r_{1}r_{2} - s_{1}s_{2}) - (r_{1} - r_{2})^{2} = 0$. Hence $Q(P) = 2|B_{\Lambda}|^{2} + 2\lambda C$, and $2\lambda = Q(I_{\Lambda}) = \Lambda/3$ by Lemma \claimref{mass.de-sitter-scale-frozen-point}(3). The identities for $C$ in \eqref{eq:gram} are the $2\times 2$ identities $\tr(h_{0}^{2}) = \tr h^{2} - \tfrac12(\tr h)^{2}$ and $\tr h^{2} = (\tr h)^{2} - 2\det h$.

(3) At $\lambda = 0$ the formula reduces to its first term, and (1) identifies $B_{0}$ with $\br{\pi^{1}\pi^{2}} = \ip{n_{\infty}}{Z^{1}\wedge Z^{2}}$.

(4) In the chart of Fact \claimref{mass.imported-flat-model-twistor-chart-infinity} the line of $x$ is the graph $L_{x} = \{\mu_{x}\pi + \pi : \pi \in \Spi\}$ of the map $\mu_{x}: \Spi \to \Som$, $(\mu_{x}\pi)^{a} = x^{aa'}\varepsilon_{b'a'}\pi^{b'}$, that is $\mu_{x} = -b_{x}$ in the notation of Lemma \claimref{mass.chart-form-spin-isomorphism} \cite[Thm.~4.2]{ButlerT9} (\claimref{correspondence.incidence-relation-chart}); indeed $(\mu_{x}e_{3} + e_{3})\wedge(\mu_{x}e_{4} + e_{4}) = n_{0} + \mu_{x}e_{3}\wedge e_{4} + e_{3}\wedge\mu_{x}e_{4} + \det(\mu_{x})\,n_{\infty} = n_{0} + x^{a4}e_{a}\wedge e_{4} + x^{a3}e_{a}\wedge e_{3} + \det(x^{aa'})\,n_{\infty} = X(x)$, since $\det\mu_{x} = \det(x^{aa'}) = -\tfrac12 Q(x)$ by Remark \ref{rem:nullity}. Two conformal spinors on $L_{x}$ are $Z^{i} = \mu_{x}\pi^{i} + \pi^{i}$, and bilinearity with antisymmetry gives $Z^{1}\wedge Z^{2} = \pi_{1}^{a'}\pi_{2}^{b'}\,(\mu_{x}e_{a'} + e_{a'})\wedge(\mu_{x}e_{b'} + e_{b'}) = \br{\pi^{1}\pi^{2}}\,(\mu_{x}e_{3} + e_{3})\wedge(\mu_{x}e_{4} + e_{4}) = \br{\pi^{1}\pi^{2}}\,X(x)$. Then $B_{\Lambda} = \br{\pi^{1}\pi^{2}}\ip{I_{\Lambda}}{X(x)} = \br{\pi^{1}\pi^{2}}\,\sigma(x)$ by Lemma \claimref{mass.de-sitter-scale-frozen-point}(2).

(5) Both terms of \eqref{eq:grip} are nonnegative for $\Lambda > 0$, and $C = 0$ iff $h_{11} = h_{22}$ and $h_{12} = 0$. For $Z^{2} = Z^{1}$: $Z\wedge Z = 0$ gives $B_{\Lambda} = 0$; $h = h_{11}\left(\begin{smallmatrix}1&1\\1&1\end{smallmatrix}\right)$ gives $C = \tfrac14\cdot 0 + h_{11}^{2}$. For the example, $\omega^{1} = e_{1}$, $\pi^{1} = e_{4}$, $\omega^{2} = \alpha e_{1}$, $\pi^{2} = \beta e_{4}$: both brackets vanish, so $B_{\Lambda} = 0$; every $t_{il} = \hp{\pi^{i}}{\omega^{l}}$ pairs a multiple of $e_{4}$ (which $\Sigma$ identifies with $e_{2}$) against a multiple of $e_{1}$ and vanishes, so $h = 0$; and $Z^{1}\wedge Z^{2} = (\beta - \alpha)\,e_{1}\wedge e_{4} \neq 0$. The script builds $P$ from the exact spin isomorphism and the moment map, on ten spinor pairs, and confirms (2) as an identity in $\lambda$ (all three coefficients of the quadratic in $\lambda$), the assembled values $Q(P_{\infty})$, $Q(P_{0})$ and $\ip{P_{\infty}}{P_{0}}$ against their closed forms with $p\cdot k$, $d$ and $\det t$ formed from the components of the proof, (3), the wedge identity and the factor $\sigma(x)$ of (4), and both parts of (5).
